\section{Related work}

\label{sec:related_work}

\noindent {\bf Fact checking and rationalized NLP models}
Fact-checking datasets include PolitiFact \cite{Vlachos2014FactCT}, Emergent \cite{ferreira-vlachos-2016-emergent}, LIAR \cite{Wang2017LiarLP}, SemEval 2017 Task 8 RumorEval \cite{derczynski-etal-2017-semeval}, Snopes \cite{Popat2017WhereTT}, CLEF-2018 CheckThat! \cite{nakov2018overview}, Verify \cite{baly-etal-2018-integrating}, Perspectrum \cite{Chen2019SeeingTF}, \fever \cite{Thorne2018FEVERAL}, and \snopes \cite{Hanselowski2019ARA}. \citet{Hanselowski2019ARA} provides a thorough review.  To our knowledge, there are no existing data sets for scientific claim verification. We refer to our task as ``claim verification'' rather than ``fact-checking'' to emphasize that our focus is to help researchers make sense of scientific findings, not to counter disinformation.

Fact-checking is one of a number of tasks where a model is required to justify a prediction via \emph{rationales} from the source document. The ERASER dataset \cite{DeYoung2019ERASERAB} provides a suite of benchmark datasets for evaluating rationalized NLP models. \ours can serve as an additional testbed for these techniques.


% \paragraph{Natural vs synthetic claims}
% We distinguish between \emph{synthetic} and \emph{natural} claims.  \fever uses synthetic claims created by annotators by mutating Wikipedia sentences selected as related evidence.  Most other prior work uses natural claims curated from fact checking sites, Twitter, debates, or news articles. The claims in \ours natural \dw{From Sewon: This seems to be an important contribution that could be emphasized in the introduction.}, since they are derived from citation sentences that occur naturally in scientific articles, and annotators to not see the evidence at time of claim writing. We discuss this claim-writing process further in \S\ref{sec:claim_writing}.

% \paragraph{Labeling claims vs claim-document pairs}
% In fact-checking, a claim is a statement of actuality whose veracity is a fixed target for investigation.  Therefore, claims can be assigned a global \emph{supported} or \emph{refuted} label. For example in \fever, the claim ``\emph{Barack Obama was the $44^{th}$ President of the United States}'' can be verified as globally \textit{supported} given sufficient evidence.

% While \ours claims are indeed factual assertions, we do not attempt to assign them global labels because the asserted ``fact'' may still be under active scientific research. Instead of labeling claims, we label claim-document pairs with \emph{support} or \emph{refute} relations.  This is similar to the task in Perspectrum \cite{Chen2019SeeingTF}, which identifies evidence-backed ``perspective'' statements as agreeing or disagreeing with an opinion-based claim, such as ``\emph{Animals should have lawful rights}.''  We discuss this claim-document labeling process further in \S\ref{sec:claim_verification}.

%%%%%%%%%%%%%%%%%%%%%%%%%%%%%%%%%%%%%%%%

\vspace{.2cm} \noindent {\bf Related scientific NLP tasks}
\label{sec:related_scinlp}
The \ours task is closely related to three other scientific NLP tasks. 
%
The \emph{citation contextualization} task \cite{Cohan2015MatchingCT,Jaidka2017TheCS} is to identify all spans in a cited document that are relevant to a particular citation in a citing document. Expert annotators achieved only 21.7\% inter-annotator agreement on the BioMedSumm dataset \cite{BiomedSumm}. 

The \emph{evidence inference} task \cite{Lehman2019InferringWM, DeYoung2020EvidenceI2}, involves predicting the effect of a medical intervention on a specified outcome. This task can be viewed as verifying a specific subset of scientific claims.
% , the evidence inference task requires the model to identify evidence justifying its label predictions. Unlike the full-sentence claims given as input to \ours, the inputs for evidence inference are individual text spans specifying an intervention, comparator, and treatment outcome.

\emph{Systematic reviews} are the ``gold standard'' for assessing the truth of scientific findings\footnote{\url{https://www.cochranelibrary.com/about/about-cochrane-reviews}}. Research on \emph{automated evidence synthesis} \cite{Marshall2019TowardSR, Beller2018MakingPW, Tsafnat2014SystematicRA, Marshall2017AutomatingBE} seeks to automate some steps of the systematic review process. Most work on evidence synthesis has focused on extracting relevant documents and phrases such as PICO snippets \cite{Nye2018ACW}. We hope that systems for claim verification will serve as components in future evidence synthesis frameworks.
